\section{Factor-model specification and support proofs}\label{app:support}
\subsection{The full conditional-variance map}\label{app:svderivation}
Put $\Gamma_{i,j}=\sum_{h=1}^i\gamma_{h,j}$, $\Gamma_{0,j}=0$, and
\[
h_j(s,T)=a_j(s)e^{-\kappa_j(T-s)}\Gamma_{\nu(T)-\nu(s),j},
\qquad \sigma_j(s,T)=h_j(s,T)\sqrt{Z_j(s)}.
\]
The HJM drift creates a random contribution to the short-rate jump. Define
\begin{align}
b_{nj}(s)&=g_{nj}(s)\int_s^{T_n}h_j(s,x)\,\dd x,&
R_{nj}(u)&=\int_u^{T_n}b_{nj}(s)e^{-\vartheta_j(s-u)}\,\dd s,\nonumber\\
\mathcal H_{nj}(u)&=g_{nj}(u)^2+2\rho_j(u)g_{nj}(u)\alpha_j(u)R_{nj}(u)
                         +\alpha_j(u)^2R_{nj}(u)^2.\label{eq:svkernel}
\end{align}
All loadings and coefficients in these expressions are deterministic and fixed.

\begin{proposition}[Conditional variance with leverage]\label{prop:svmap}
Under these hypotheses the actual conditional short-rate jump variances satisfy
the affine identity \eqref{eq:svmap},
where
\begin{align*}
\mathsf L_{nj}(t)&=\int_t^{T_n}\mathcal H_{nj}(u)e^{-\vartheta_j(u-t)}\,\dd u,\\
\boldsymbol c_n(t)&=\sum_j\int_t^{T_n}\mathcal H_{nj}(u)(1-e^{-\vartheta_j(u-t)})\,\dd u.
\end{align*}
Both are nonnegative. For fixed parameters the variance vector therefore lies in a translated cone of dimension at most the number of factors.
\end{proposition}
\begin{proof}
The maturity discontinuity of the differentiated HJM equation gives
\[
\Delta r_{T_n}=\Delta f_0(T_n)+\sum_j\left[\int_0^{T_n}b_{nj}(s)Z_j(s)\,\dd s
                  +\int_0^{T_n}g_{nj}(s)\sqrt{Z_j(s)}\,\dd W_j(s)\right].
\]
The initial-curve jump $\Delta f_0(T_n)$ is deterministic and drops out of the conditional variance. Since $R'_{nj}=\vartheta_jR_{nj}-b_{nj}$ and $R_{nj}(T_n)=0$, integration by parts replaces the random part of the future drift integral by $\int_t^{T_n}\alpha_jR_{nj}\sqrt{Z_j}\,\dd U_j$. Conditional isometry yields
\[
\boldsymbol V_n(t)=\sum_j\int_t^{T_n}\mathcal H_{nj}(u)
                      \E[Z_j(u)\mid\F_t]\,\dd u.
\]
Use $\E[Z_j(u)\mid\F_t]=1+(Z_j(t)-1)e^{-\vartheta_j(u-t)}$. Finally,
$\mathcal H_{nj}=(g_{nj}+\rho_j\alpha_jR_{nj})^2+(1-\rho_j^2)(\alpha_jR_{nj})^2\geq0$.
\end{proof}
Keeping only the integral of $g_{nj}^2\E[Z_j]$ omits both the random drift variance and its leverage covariance. Equation~\eqref{eq:svkernel} supplies the correction needed before drawing conclusions about the cross-meeting restrictions of a stochastic-volatility calibration.


\subsection{Mapping to the published scheduled-rate specification}\label{app:bgsmapping}
The mapping uses equations (3), (35)--(41), and (84) of \citet{brace2024sofr}, with their calibration Table~2. Superscript $\mathrm B$ below distinguishes a source symbol when its meaning differs from ours. Their meeting count is
$A^{\mathrm B}_{s,T}=\#\{n:s<T_n\leq T\}=\nu(T)-\nu(s)$.
Equations (36)--(38) give the forward-rate volatility
\begin{align*}
\sigma_j^{\mathrm B}(s,T)
 &=\chi_j^{\mathrm B}(s)\phi_j^{\mathrm B}(T)
       \sum_{i=1}^{A^{\mathrm B}_{s,T}}\gamma_{i,j},\\
\phi_j^{\mathrm B}(T)&=\exp\left(-\int_0^T\lambda_j^{\mathrm B}(u)\,\dd u\right),\\
\chi_j^{\mathrm B}(s)&=\sigma_j^{\mathrm B}(s)\sqrt{v_j^{\mathrm B}(s)}
          \exp\left(\int_0^s\lambda_j^{\mathrm B}(u)\,\dd u\right).
\end{align*}
In the last line $\sigma_j^{\mathrm B}(s)$ is their scalar volatility scale, not their two-argument forward loading. In Table~2 these scales and the forward decay rates are constant in both columns; only the $\alpha_j$ coefficients have multiple time pieces. Thus their equation (84) becomes exactly
\[
\sigma_j^{\mathrm B}(s,T)
=a_j(s)e^{-\kappa_j(T-s)}\Gamma_{\nu(T)-\nu(s),j}\sqrt{Z_j(s)}
=h_j(s,T)\sqrt{Z_j(s)},
\]
with the following correspondence.

\begin{center}\small
\begin{tabular}{lll}
\toprule
This paper & Source symbol & Role\\
\midrule
$T_n$ & $x_n$ & Scheduled meeting date\\
$\gamma_{i,j}$ & $\gamma_{i,j}$ & Ordinal loading, unchanged\\
$\Gamma_{m,j}$ & $\sum_{i=1}^{m}\gamma_{i,j}$ & Cumulative ordinal loading\\
$a_j(s)$ & $\sigma_j^{\mathrm B}(s)$ & Deterministic scale\\
$\kappa_j$ & Constant $\lambda_j^{\mathrm B}$ & Forward-loading decay\\
$Z_j(s)$ & $v_j^{\mathrm B}(s)$ & Variance state, initially one\\
$\vartheta_j$ & Constant $\theta_j^{\mathrm B}$ & Variance-state mean reversion\\
$\alpha_j(s),\rho_j$ & $\alpha_j^{\mathrm B}(s),\rho_j^{\mathrm B}$ & Volatility of volatility and leverage\\
\bottomrule
\end{tabular}
\end{center}
Their equations (39)--(41) give \eqref{eq:svstate} and the within-factor correlation. Their Section~3.6 and footnote~29 specify independent driving factors. The general source permits time-dependent $\lambda_j$ and $\theta_j$; our constant-decay and constant-speed assumptions specialize it exactly to the two printed calibrations.

At a meeting maturity, the sum of ordinal loadings acquires exactly one term. Therefore, for $s<T_n$,
\[
h_j(s,T_n)-h_j(s,T_n-)
=a_j(s)e^{-\kappa_j(T_n-s)}\gamma_{n-\nu(s),j}=g_{nj}(s).
\]
Their HJM drift, equation (35), has factor contribution
$h_j(s,T)Z_j(s)\int_s^T h_j(s,x)\,\dd x$.
The maturity integral is continuous at $T_n$, so its jump coefficient is exactly
$b_{nj}(s)=g_{nj}(s)\int_s^{T_n}h_j(s,x)\,\dd x$.
This proves the loading and drift correspondence used in Proposition~\ref{prop:svmap}; it is not an identification of their option prices with the announcement model's Gaussian cash formula.

For direct comparison, the complete parameter columns are:
\begin{center}\small
\begin{tabular}{lll}
\toprule
Parameter & Constant column & Time-dependent column\\
\midrule
$(a_1,a_2,a_3)$ & $(0.0081,0.006,0.0041)$ & $(0.00663,0.00587,0.00447)$\\
$(\kappa_1,\kappa_2,\kappa_3)$ & $(0.01,0,0.16)$ & $(0.02,0.004,0.35)$\\
$\alpha_1$ & $1.57$ & $(3.142,1.35,3.2,0.86)$\\
$\alpha_2$ & $0.82$ & $(0.76,0.66,0.6,0.22)$\\
$\alpha_3$ & $0$ & $(3.0,0.6,0.5,4.1)$\\
$(\rho_1,\rho_2,\rho_3)$ & $(-0.2,0,0)$ & $(-0.14,-0.025,-0.83)$\\
$(\vartheta_1,\vartheta_2,\vartheta_3)$ & $(0,0,0)$ & $(0.1,0,11.0)$\\
\bottomrule
\end{tabular}
\end{center}
The source's Section~4.2 places the four $\alpha_j$ values on successive option-expiry intervals. Each is positive in the time-dependent column, so their exact endpoints are unnecessary for the active-factor count on that calibrated horizon. They and the ordinal loadings would be necessary to compute the numerical columns of $\mathsf L$.

In the constant column, $Z_3(t)=1$ and $\boldsymbol c(t)=0$. Consequently
\[
\operatorname{supp}\boldsymbol V(t)=\mathsf L_{\cdot3}(t)
 +\operatorname{cone}\{\mathsf L_{\cdot1}(t),\mathsf L_{\cdot2}(t)\},
\]
of affine dimension at most two. In the time-dependent column all three states are active, so the support is $\boldsymbol c(t)+\mathsf L(t)\R_+^3$, with dimension $\rank\mathsf L(t)\leq3$. Three active states do not imply three linearly independent loading columns. With $m$ future meetings, the two specifications therefore have at least $m-2$ and $m-3$ affine equalities, respectively. These statements concern each fixed parameter specification, not a panel pooled over daily recalibrations.

\subsection{Proof of the support theorem}
We prove Theorem~\ref{thm:svsupport} under the solution and integrability hypotheses in Section~\ref{sec:restrictions}. On one coefficient piece of length $h$, starting from $x\geq0$, with $\alpha>0$ and $\vartheta>0$, the square-root transition has Laplace transform
\begin{equation}\label{eq:cirtransform}
\E[e^{-uZ(s+h)}\mid\F_s]
=(1+c_hu)^{-2\vartheta/\alpha^2}
 \exp\left(-\frac{u e^{-\vartheta h}x}{1+c_hu}\right),
\qquad c_h=\frac{\alpha^2}{2\vartheta}(1-e^{-\vartheta h}).
\end{equation}
This is the classical scaled noncentral chi-square transition; see \citet{malham2008square}, Proposition~1. For $\vartheta=0$ it becomes
$\exp[-ux/(1+\alpha^2hu/2)]$, and for $\alpha=0$ the transition is the deterministic flow $1+(x-1)e^{-\vartheta h}$.

These formulas apply conditionally to the assumed solutions. One can obtain them directly by applying Itô's formula to the exponential-affine function with terminal value $e^{-uZ}$: its backward coefficients cancel the drift. Stop on bounded state intervals to make the stochastic integral a true martingale, and use bounded convergence for the exponential, which lies in $[0,1]$. Applying the same argument jointly, the independent factor drivers have zero cross variations, and the joint transform is the product of the coordinate transforms. Iteration across the finite common partition proves the conditional product transition law, without assuming an additional independence property of the chosen solutions.

For $\vartheta>0$, the law in \eqref{eq:cirtransform} has positive density on $(0,\infty)$, support $[0,\infty)$, and no atom at zero. This also holds when $x=0$. For $\vartheta=0$ and $x>0$, put $c=\alpha^2h/2$. The transform is that of a sum of a Poisson number of independent exponential variables, with Poisson mean $x/c$ and exponential mean $c$. Its support is $[0,\infty)$ and its zero mass is $e^{-x/c}$; at $x=0$ it is absorbed. A deterministic final piece maps a half-line to $[\ell_j,\infty)$ and preserves any mass at its lower endpoint. Starting from one, any positive-length stochastic piece therefore activates that coordinate, and each active coordinate has strictly positive variance. These facts prove
\[
\operatorname{supp}Z(t)=
 \prod_{j\notin J}\{1\}\ \times\prod_{j\in J}[\ell_j,\infty),
\]
with the product interpreted in coordinate order.

Apply the continuous affine map $z\mapsto\boldsymbol c+\mathsf Lz$. Its image is the finitely generated translated cone in \eqref{eq:support}, which is closed. The support of the image law is the closure of the image of the support, so here it equals that cone. Its affine hull is the stated translate of $\operatorname{range}\mathsf L_J$. Coordinate independence gives the covariance formula. For any vector $\omega$,
\[
\omega^\top\Cov(\boldsymbol V(t))\omega
 =\sum_{j\in J}\Var(Z_j(t))(\omega^\top\mathsf L_{\cdot j})^2.
\]
All coefficients are positive, hence the covariance kernel is $\ker\mathsf L_J^\top$. This proves both the rank and the exhaustive affine-equality statements.

For completeness, at $0\leq s<t$, the regular conditional version uses
\[
J_s=\{j:\alpha_j>0\text{ on a positive-length piece of }(s,t],
                   \ \vartheta_j>0\text{ or }Z_j(s)>0\}.
\]
Compute the lower endpoints from the state $Z(s)$ and the intervening pieces; an inactive coordinate follows its deterministic flow. The same support formula holds with these endpoints and $J_s$, almost surely. Its conditional covariance is the sum of the column outer products weighted by the conditional coordinate variances. In particular, a zero-speed coordinate absorbed at zero supplies no later random direction, whereas positive mean reversion allows a zero coordinate to become active on a later stochastic piece.

Since the columns of $\mathsf L$ are nonnegative, reaching the cone vertex requires every active coordinate with a nonzero column to reach its lower endpoint. The vertex mass is the product of those endpoint masses. Any such coordinate with positive mean reversion makes the product zero. For zero speeds and constant positive $\alpha_j$, starting from one, the vertex mass is
\[
\exp\left(-\sum_{\substack{j\in J\\\mathsf L_{\cdot j}\ne0}}
                            \frac{2}{\alpha_j^2t}\right).
\]
This distinguishes a support boundary, which can always be approached, from an atom at that boundary.
